% ============================================================
% CAPITOLO 16 - SCHEMI A MODULI DI DUE ACCORDI
% ============================================================
\chapter{Schemi a moduli di due accordi}
\label{chap:moduli}

Una progressione ripete un \textbf{modulo di due accordi} su gradi diversi. Il nome dello schema indica il movimento del basso nel modulo e il verso in cui il modulo si sposta: \textbf{I--IV} significa che il basso va dal primo accordo a quello una quarta sopra (Do--Fa), \textbf{I--V} che va una quinta sopra (Do--Sol); \textbf{discendente} o \textbf{ascendente} dice che ogni nuovo modulo parte un grado sotto o sopra il precedente.

Si distinguono due famiglie:
\begin{itemize}
	\item la progressione \textbf{non modulante} (diatonica), che usa solo gli accordi di Do maggiore e torna alla tonica;
	\item la progressione \textbf{modulante} (reale), che ripete il modulo identico trasposto di un tono per volta; dopo ogni modulo si indica il grado di arrivo rispetto alla tonalità di partenza.
\end{itemize}

In questo capitolo gli schemi sono realizzati con le sole triadi. Le versioni con settime sono nel Capitolo \ref{chap:progressioni_settime}.

\begin{regola}[Condotta nelle progressioni a moduli]
	\label{reg:condotta_moduli}
	\begin{itemize}
		\item Niente quinte e ottave parallele, né intervalli melodici eccedenti o diminuiti nelle voci superiori.
		\item La sensibile non si raddoppia.
		\item Le triadi sono complete; si accetta la triade incompleta (fondamentale triplicata) solo quando serve a una buona condotta.
		\item Nell'ultimo accordo la sensibile del V sale alla tonica.
		\item Le voci restano nella zona preferibile e ben separate (Definizione \ref{def:estensione_voci}).
	\end{itemize}
\end{regola}

\section{Moduli I--IV}
\label{sec:moduli_i4}

\subsection{Progressione I--IV discendente}
\label{subsec:prog_i4d}

Modulo: il primo accordo, poi l'accordo sul IV grado; il modulo si ripete un grado sotto per volta.

\begin{esempio}[I--IV discendente in Do maggiore: non modulante]
	\label{ex:prog_i4d_dia}
	\musica{es_cap17_i4d_dia_T}
\end{esempio}

\begin{esempio}[I--IV discendente: modulante, sequenza reale di tre moduli]
	\label{ex:prog_i4d_rea}
	\musica{es_cap17_i4d_rea_T}
\end{esempio}

\subsection{Progressione I--IV ascendente}
\label{subsec:prog_i4a}

Modulo: il primo accordo, poi l'accordo sul IV grado; il modulo si ripete un grado sopra per volta.

\begin{esempio}[I--IV ascendente in Do maggiore: non modulante]
	\label{ex:prog_i4a_dia}
	\musica{es_cap17_i4a_dia_Ta}
	\musica{es_cap17_i4a_dia_Tb}
\end{esempio}

\begin{esempio}[I--IV ascendente: modulante, sequenza reale di tre moduli]
	\label{ex:prog_i4a_rea}
	\musica{es_cap17_i4a_rea_T}
\end{esempio}

\section{Moduli I--V}
\label{sec:moduli_i5}

\subsection{Progressione I--V ascendente}
\label{subsec:prog_i5a}

Modulo: il primo accordo, poi l'accordo sul V grado; il modulo si ripete un grado sopra per volta.

\begin{esempio}[I--V ascendente in Do maggiore: non modulante]
	\label{ex:prog_i5a_dia}
	\musica{es_cap17_i5a_dia_T}
\end{esempio}

\begin{esempio}[I--V ascendente: modulante, sequenza reale di tre moduli]
	\label{ex:prog_i5a_rea}
	\musica{es_cap17_i5a_rea_T}
\end{esempio}

\subsection{Progressione I--V discendente}
\label{subsec:prog_i5d}

Modulo: il primo accordo, poi l'accordo sul V grado; il modulo si ripete un grado sotto per volta.

\begin{esempio}[I--V discendente in Do maggiore: non modulante]
	\label{ex:prog_i5d_dia}
	\musica{es_cap17_i5d_dia_Ta}
	\musica{es_cap17_i5d_dia_Tb}
\end{esempio}

\begin{esempio}[I--V discendente: modulante, sequenza reale di tre moduli]
	\label{ex:prog_i5d_rea}
	\musica{es_cap17_i5d_rea_T}
\end{esempio}

\section{Progressioni imitate}
\label{sec:progressioni_imitate}

Nella progressione \textbf{imitata} il primo modulo è scritto liberamente e ogni modulo successivo lo ripete \textbf{identico}, un grado sopra o sotto: ogni voce compie lo stesso movimento, per lo stesso numero di gradi, del modulo precedente. Gli accordi sono quelli diatonici di Do maggiore e l'ultimo non è risolto, perché la progressione termina con il modulo. In una ripetizione fedele del modello può comparire il raddoppio della sensibile (nel VII$^\circ$), perché in una sequenza l'accordo non ha funzione di dominante.

\begin{esempio}[Progressione imitata I--IV discendente]
	\label{ex:imitata_i4d}
	\musica{es_cap17_i4d_imi_T}
\end{esempio}

\begin{esempio}[Progressione imitata I--IV ascendente]
	\label{ex:imitata_i4a}
	\musica{es_cap17_i4a_imi_Ta}
	\musica{es_cap17_i4a_imi_Tb}
\end{esempio}

\begin{esempio}[Progressione imitata I--V ascendente]
	\label{ex:imitata_i5a}
	\musica{es_cap17_i5a_imi_T}
\end{esempio}
